\documentclass[11pt,reqno]{amsart}
\usepackage[T1]{fontenc}
\usepackage{lmodern,amsmath,amssymb,amsthm,mathtools,microtype,mathrsfs}
\usepackage[margin=1.1in]{geometry}
\usepackage{longtable,booktabs,array}
\usepackage{xr-hyper}
\usepackage[hidelinks]{hyperref}
\allowdisplaybreaks[2]
\setlength{\emergencystretch}{2em}
\numberwithin{equation}{section}
\newtheorem{theorem}{Theorem}[section]
\newtheorem{lemma}[theorem]{Lemma}
\newtheorem{proposition}[theorem]{Proposition}
\newtheorem{corollary}[theorem]{Corollary}
\theoremstyle{definition}\newtheorem{definition}[theorem]{Definition}
\theoremstyle{remark}\newtheorem{remark}[theorem]{Remark}
\newcommand{\R}{\mathbb R}\newcommand{\Z}{\mathbb Z}
\newcommand{\E}{\mathbb E}\newcommand{\dd}{\,\mathrm d}
\newcommand{\one}{\mathbf1}\newcommand{\csch}{\operatorname{csch}}
\newcommand{\arcosh}{\operatorname{arcosh}}\newcommand{\tr}{\operatorname{tr}}
\newcommand{\diag}{\operatorname{diag}}\newcommand{\supp}{\operatorname{supp}}
\newcommand{\cC}{\mathcal C}\newcommand{\cD}{\mathcal D}
\newcommand{\cP}{\mathcal P}\newcommand{\id}{\mathrm{id}}
\raggedbottom
\hypersetup{pdftitle={Action rigidity with selective detection},pdfauthor={Qian Qi},pdfsubject={A2 revision 83: sharp clocks and global quotient reconstruction}}
\title[Action rigidity with selective detection]{Action rigidity with selective detection}
\author{Qian Qi}
\date{September 18, 2026}
\subjclass[2020]{41A05, 37J10, 53D10, 62H12}
\keywords{Generating actions, Chebyshev systems, projective observations,
 selective detection, finite mixtures, covering spaces, return maps}
\begin{document}
\begin{abstract}
We recover absolute generating actions from projective endpoint observations
with unknown selective detection. If the logarithm of the branchwise detector
belongs to the polynomials of degree at most $q$, exactly $q+3$ deadlines
suffice in the worst case. Neither the number of visible branches nor the
two independent readout channels is known. The proof reduces the action
parameters to the endpoints of an oriented interval determined by two
weighted moments. More generally, for a two-dimensional annihilator
containing a positive Cauchy kernel, scalar identification is equivalent
to a strict two-function Chebyshev condition. We construct a global
$C^m$ inverse near the exact data, modulo permutation and the common
detector gauge, by quantitative synchronization on the component covering.
This includes nontrivial monodromy and perturbations outside the model.
The recovered actions determine exact return lifts on covered phase domains.
A localized Hamiltonian ambiguity identifies an obstruction to removing
phase coverage in the smooth category, while a nonconvex three-sheet
return system realizes the positive theorem through an action crossing.
\end{abstract}
\maketitle
\setcounter{tocdepth}{1}
\tableofcontents
